\chapter{Explainability and pathway prioritisation}\label{ch:explainability}

\section{Gene-set importance}\label{sec:importance}

All 50 fits of M2 completed and were analysed. Ranked by mean held-out permutation importance
(Equation~\ref{eq:importance}; Figure~\ref{fig:importance}a; Table~\ref{tab:s2}), the leading
gene set is \textbf{HALLMARK\_P53\_PATHWAY}, with mean importance 0.00343 mSF2 units,
importance above zero in 90\% of fits, sign consistency 0.92 and a predominantly negative
coefficient. It is followed by COMPLEMENT (0.00118), OXIDATIVE\_PHOSPHORYLATION (0.00107),
UV\_RESPONSE\_UP (0.00100), MYC\_TARGETS\_V2, KRAS\_SIGNALING\_UP and UV\_RESPONSE\_DN.

Three observations matter more than the ranking itself.
\begin{itemize}
  \item \textbf{The magnitudes are tiny.} The largest mean importance, 0.0034 mSF2 units, is
    of the same order as the entire predictive gain, so the ranking describes a very small
    effect.
  \item \textbf{Stability is modest.} Across the 50 fits, the mean Kendall $\tau$ between
    importance rankings is 0.225, the mean rank-biased overlap 0.518 and the mean top-10
    Jaccard index 0.391: the same gene sets do not reliably appear at the top of every fit.
  \item \textbf{Correlated sets share signal.} Several ranked sets are correlated with
    another set at $\lvert\rho\rvert\ge0.80$ (Figure~\ref{fig:s-correlation}), including
    COMPLEMENT with APOPTOSIS and MYC\_TARGETS\_V2 with MYC\_TARGETS\_V1, so importance may
    be shared or exchanged between them rather than attributable to one.
\end{itemize}

Equally important is what the ranking does \emph{not} contain. \textbf{DNA repair, G2M
checkpoint and hypoxia, the pathways a radiobiologist would expect, were not prioritised}:
G2M checkpoint ranked 13th with importance above zero in only 8\% of fits, hypoxia 10th with
10\%, and DNA repair 47th. Reporting this guards against presenting a familiar
radiation-response narrative that the data do not support.

\begin{figure}[tbp]
  \centering
  \includegraphics[width=\linewidth,height=\ReportFigureMaxHeight,keepaspectratio]{fig_06_gene_set_importance}
  \caption[Exploratory importance of the Hallmark gene-set scores in M2]{Exploratory
    importance of the Hallmark gene-set scores in the expression model M2. (a)~The fifteen
    highest-ranked gene sets. Grey points are the grouped permutation importance in each of
    the fifty model fits (ten repeats by five folds), computed on held-out rows only as the
    increase in component-weighted MAE when a score is permuted across components; diamonds
    are the means over fits, coloured by the prespecified tier. The spread across fits is
    split-to-split variability, not a confidence interval. (b)~The two tier-C criteria for
    all fifty gene sets: the share of fits with importance above zero and the share of fits
    in which the set received a non-zero coefficient; the shaded corner satisfies both
    thresholds (dashed lines at 0.50). Tiers A and B were unreachable by design. Shares are
    proportions of overlapping cross-validation fits, not probabilities. The rankings
    describe one fitted model family on one dataset; they are not validated biomarkers,
    causal mechanisms, radiation-response pathways or molecular effect sizes.}
  \label{fig:importance}
\end{figure}

\section{The pathway shortlist}\label{sec:pathway-shortlist}

Seven of the 50 gene sets reached \textbf{tier~C} and 43 were not prioritised; \textbf{no set
reached tier A or B}, because both require a formally demonstrated prediction gain and at
least 1,000 negative-control draws, and the run has neither (Figure~\ref{fig:importance}b).
The gene-set tests record zero completed draws of both the permuted-outcome and the
random-gene-set nulls, so attribution is not permitted and the stability gate cannot be
evaluated.

The shortlist has a precise meaning. Its members are the gene-set scores whose permutation
importance was positive in at least half of the cross-validation fits and which received a
non-zero coefficient in at least half, within one model family on one dataset. It is not a
gene-level finding, a validated biomarker, a causal mechanism, a radiation-response pathway
or a molecular effect size, and it is not evidence that modulating any of these pathways
would change radiosensitivity.

\section{Cell lines the model reproduces least well}\label{sec:discrepancies}

Of the 241 measured lines, 12 reached discrepancy tier~A, nine tier~B and 27 tier~C; 193 were
not prioritised (Figure~\ref{fig:discrepancies}; Table~\ref{tab:s3}). The highest-ranked
discrepancies are dominated by lines with extreme observed values --- HT-144 (observed 0.08,
predicted 0.59 under M2), KNS-42 (0.125 against 0.60), HEC-1-B (0.97 against 0.56) and
D283~Med (0.18 against 0.58) --- which is the shrinkage of Figure~\ref{fig:prediction} seen one line at a time
rather than a set of biologically unusual cell lines.

The triage separates possible explanations. Several tier A and B lines carry technical flags
(small tissue group, related lines in the cohort, a Cellosaurus problem flag on a relative)
and are routed to technical review first; the remainder are labelled as unexplained
discrepancies requiring follow-up. Because the tier thresholds are cohort percentiles, a
roughly fixed fraction of lines is always flagged: a flag means \enquote{relatively large in
this cohort}, not \enquote{abnormal}.

Of the two context lines from the motivating study, \textbf{MCF-7 was not prioritised}: its
held-out errors lie within the bulk of the cohort (observed 0.52 against 0.574 under M1 and
0.534 under M2). \textbf{DU145 reached tier~C} with four technical flags --- stranded RNA
profile, small tissue group, unstable prediction across repeats and strong model
disagreement --- and was observed at 0.71 against predictions of 0.580 under M1 and 0.650
under M2. Both tiers were assigned by the same rule as for every other line. No individual
prediction error is evidence of radioprotection, radiosensitisation or any causal mechanism,
and lines selected for extreme residuals are expected to agree better on re-measurement even
if nothing is wrong, through regression to the mean.

\begin{figure}[tbp]
  \centering
  \includegraphics[width=\linewidth,height=\ReportFigureMaxHeight,keepaspectratio]{fig_07_cell_line_discrepancies}
  \caption[Cell lines whose radiosensitivity the model reproduces least well]{Cell lines whose
    measured radiosensitivity the model reproduces least well. (a)~Residual (observed minus
    predicted mSF2 under M2, predictions averaged over repeats) against the observed value
    for every measured line, coloured by prespecified discrepancy tier; triangles mark lines
    with at least one technical flag. The strong dependence of the residual on the observed
    value is shrinkage: because the models predict close to the cohort mean, lines with
    extreme measured values necessarily have the largest residuals. Circled points are MCF-7
    and DU145, labelled for their relevance to the motivating study; their tiers are not
    altered. (b)~Flags recorded for the tier A and B lines, in rank order, separating flags
    about prediction behaviour (instability across repeats, disagreement between M1 and M2)
    from flags about the data or cell-line identity (stranded RNA profile, related lines in
    the cohort, a Cellosaurus problem flag on a relative, small tissue group). These are
    measured lines used in model development, not independent validation samples; a large
    residual is a signal to re-check identity, provenance and the original measurement, not
    evidence of unusual radiobiology.}
  \label{fig:discrepancies}
\end{figure}
